\documentclass[11pt,a4paper]{article}
\usepackage[utf8]{inputenc}
\usepackage{amsmath,amssymb,amsthm}
\usepackage{algorithm}
\usepackage{algpseudocode}
\usepackage{booktabs}
\usepackage{hyperref}
\usepackage{geometry}
\geometry{margin=1in}

\title{Gated Recurrent Unit (GRU): Two-Gate Parameter Efficiency, Linear State Interpolation, and Dynamic Temporal Reset}
\author{Automorphic Intelligence Architecture Specifications}
\date{\today}

\newtheorem{theorem}{Theorem}
\newtheorem{lemma}[theorem]{Lemma}
\newtheorem{proposition}[theorem]{Proposition}
\theoremstyle{definition}
\newtheorem{definition}{Definition}
\newtheorem{remark}{Remark}

\begin{document}

\maketitle

\begin{abstract}
While Long Short-Term Memory (LSTM) architectures resolve gradient vanishing via dedicated cell states and three gating mechanisms, their parameter footprint and computational overhead remain substantial. The Gated Recurrent Unit (GRU) streamlines recurrent gating dynamics into two compact mechanisms—the update gate and the reset gate—completely eliminating the separate internal cell memory state while maintaining comparable long-range temporal retention. This document provides the formal mathematical specification of GRU transitions, derives gradient highway dynamics under update gating, provides algorithmic pseudocode, and delivers a verified step-by-step numerical example.
\end{abstract}

\section{Notation and Mathematical Preliminaries}

Let $\mathbf{x}_t \in \mathbb{R}^{d_x}$ denote the input feature vector at sequence step $t$, and $\mathbf{h}_t \in \mathbb{R}^{d_h}$ denote the recurrent hidden state vector.

\begin{table}[h]
\centering
\caption{Mathematical Notation for Gated Recurrent Unit}
\begin{tabular}{lll}
\toprule
\textbf{Symbol} & \textbf{Domain} & \textbf{Semantic Description} \\
\midrule
$\mathbf{z}_t$ & $(0, 1)^{d_h}$ & Update gate vector (controls balance of old vs candidate state) \\
$\mathbf{r}_t$ & $(0, 1)^{d_h}$ & Reset gate vector (controls access to prior history) \\
$\tilde{\mathbf{h}}_t$ & $(-1, 1)^{d_h}$ & Candidate hidden state vector \\
$W_{\{z,r,h\}}$ & $\mathbb{R}^{d_h \times d_x}$ & Input projection weight matrices \\
$U_{\{z,r,h\}}$ & $\mathbb{R}^{d_h \times d_h}$ & Recurrent hidden projection weight matrices \\
$\mathbf{b}_{\{z,r,h\}}$ & $\mathbb{R}^{d_h}$ & Bias parameter vectors \\
$\sigma(\cdot)$ & $\mathbb{R} \to (0, 1)$ & Standard logistic sigmoid function \\
$\odot$ & Operator & Hadamard elementwise product \\
\bottomrule
\end{tabular}
\end{table}

\section{Problem Definition and Core Concepts}

\subsection{3-Level Explanation}
\begin{enumerate}
    \item \textbf{Intuitive (High School level):} An LSTM has three guards and a separate notebook. A GRU is lighter: it keeps its memory directly in its head and uses only two guards. The Reset Guard decides whether to wipe out past memories, and the Update Guard decides how much new information replaces the old memory.
    \item \textbf{Engineering Level:} A GRU eliminates the separate cell state $\mathbf{c}_t$, cutting parameters by $25\%$ compared to an LSTM ($3(d_h d_x + d_h^2)$ vs $4(d_h d_x + d_h^2)$). Linear state interpolation $\mathbf{h}_t = (1 - \mathbf{z}_t) \odot \mathbf{h}_{t-1} + \mathbf{z}_t \odot \tilde{\mathbf{h}}_t$ maintains additive gradient highways.
    \item \textbf{Mathematical Level:} Given sequence $\{\mathbf{x}_t\}$, the GRU state evolution equations are defined by:
    \begin{align}
    \mathbf{z}_t &= \sigma(W_z \mathbf{x}_t + U_z \mathbf{h}_{t-1} + \mathbf{b}_z) \\
    \mathbf{r}_t &= \sigma(W_r \mathbf{x}_t + U_r \mathbf{h}_{t-1} + \mathbf{b}_r) \\
    \tilde{\mathbf{h}}_t &= \tanh(W_h \mathbf{x}_t + U_h (\mathbf{r}_t \odot \mathbf{h}_{t-1}) + \mathbf{b}_h) \\
    \mathbf{h}_t &= (1 - \mathbf{z}_t) \odot \mathbf{h}_{t-1} + \mathbf{z}_t \odot \tilde{\mathbf{h}}_t
    \end{align}
\end{enumerate}

\subsection{5-Step Algorithmic Framework}
\begin{enumerate}
    \item \textbf{Update Gate Evaluation:} Compute $\mathbf{z}_t = \sigma(W_z \mathbf{x}_t + U_z \mathbf{h}_{t-1} + \mathbf{b}_z)$.
    \item \textbf{Reset Gate Evaluation:} Compute $\mathbf{r}_t = \sigma(W_r \mathbf{x}_t + U_r \mathbf{h}_{t-1} + \mathbf{b}_r)$.
    \item \textbf{Gated History Masking:} Multiply previous state by reset gate $\mathbf{h}_{\text{reset}} = \mathbf{r}_t \odot \mathbf{h}_{t-1}$.
    \item \textbf{Candidate State Synthesis:} Compute non-linear candidate $\tilde{\mathbf{h}}_t = \tanh(W_h \mathbf{x}_t + U_h \mathbf{h}_{\text{reset}} + \mathbf{b}_h)$.
    \item \textbf{Convex State Interpolation:} Blend prior state and candidate: $\mathbf{h}_t = (1 - \mathbf{z}_t) \odot \mathbf{h}_{t-1} + \mathbf{z}_t \odot \tilde{\mathbf{h}}_t$.
\end{enumerate}

\section{Algorithm Specification}

\begin{algorithm}[H]
\caption{Gated Recurrent Unit (GRU) Step and Unrolling}
\label{alg:gru}
\begin{algorithmic}[1]
\Require Sequence $\{\mathbf{x}_t\}_{t=1}^T$, initial state $\mathbf{h}_0 = \mathbf{0}$, weights $W_z, U_z, W_r, U_r, W_h, U_h$, biases $\mathbf{b}_z, \mathbf{b}_r, \mathbf{b}_h$.
\Ensure Hidden state trajectory $\{\mathbf{h}_t\}_{t=1}^T$.
\For{$t = 1$ \textbf{to} $T$}
    \State $\mathbf{z}_t \gets \sigma(W_z \mathbf{x}_t + U_z \mathbf{h}_{t-1} + \mathbf{b}_z)$
    \State $\mathbf{r}_t \gets \sigma(W_r \mathbf{x}_t + U_r \mathbf{h}_{t-1} + \mathbf{b}_r)$
    \State $\tilde{\mathbf{h}}_t \gets \tanh(W_h \mathbf{x}_t + U_h (\mathbf{r}_t \odot \mathbf{h}_{t-1}) + \mathbf{b}_h)$
    \State $\mathbf{h}_t \gets (\mathbf{1} - \mathbf{z}_t) \odot \mathbf{h}_{t-1} + \mathbf{z}_t \odot \tilde{\mathbf{h}}_t$
\EndFor
\State \Return $\{\mathbf{h}_t\}_{t=1}^T$
\end{algorithmic}
\end{algorithm}

\section{Theoretical Guarantees and Gradient Highway Properties}

\begin{theorem}[Gradient Highway under Zero Update Gating]
\label{thm:gru_highway}
When the update gate is closed ($\mathbf{z}_k \to \mathbf{0}$) across a span of temporal steps $k \in [t+1, T]$, the Jacobian derivative of the hidden state satisfies:
\begin{equation}
\frac{\partial \mathbf{h}_T}{\partial \mathbf{h}_t} = \prod_{k=t+1}^T \text{diag}(\mathbf{1} - \mathbf{z}_k) = I_{d_h}
\end{equation}
preserving the error signal indefinitely without temporal attenuation.
\end{theorem}

\begin{proof}
Differentiating $\mathbf{h}_k = (1 - \mathbf{z}_k) \odot \mathbf{h}_{k-1} + \mathbf{z}_k \odot \tilde{\mathbf{h}}_k$ when $\mathbf{z}_k \to \mathbf{0}$ yields $\frac{\partial \mathbf{h}_k}{\partial \mathbf{h}_{k-1}} = \text{diag}(1 - \mathbf{z}_k) = I$. Taking the product over $T - t$ steps gives the identity operator.
\end{proof}

\section{Complexity Analysis}

\begin{itemize}
    \item \textbf{Time Complexity:} $\mathcal{O}(3 \cdot (d_h d_x + d_h^2))$ FLOPs per time step ($25\%$ fewer FLOPs than standard LSTM).
    \item \textbf{Space Complexity:} $\mathcal{O}(d_h)$ active memory for the single state vector $\mathbf{h}_t$.
\end{itemize}

\section{Exactness and Error Analysis}

The convex linear interpolation $\mathbf{h}_t = (1 - \mathbf{z}_t)\mathbf{h}_{t-1} + \mathbf{z}_t \tilde{\mathbf{h}}_t$ guarantees $\|\mathbf{h}_t\|_\infty \le \max(\|\mathbf{h}_{t-1}\|_\infty, \|\tilde{\mathbf{h}}_t\|_\infty) \le 1$, ensuring bounded state dynamics.

\section{Worked Numerical Example}

Let $d_x = 1, d_h = 1, x_1 = 1.0, h_0 = 0.0$.
Parameters:
$W_z = [0.0], U_z = [0.0], b_z = [0.0] \implies z_1 = \sigma(0) = 0.5$.
$W_r = [0.0], U_r = [0.0], b_r = [0.0] \implies r_1 = \sigma(0) = 0.5$.
$W_h = [1.0], U_h = [1.0], b_h = [0.0]$.
\begin{itemize}
    \item Reset state: $r_1 \times h_0 = 0.5 \times 0.0 = 0.0$.
    \item Candidate: $\tilde{h}_1 = \tanh(1.0 \times 1.0 + 1.0 \times 0.0) = \tanh(1.0) \approx 0.7616$.
    \item State update: $h_1 = (1 - 0.5)(0.0) + 0.5(0.7616) = 0.3808$.
\end{itemize}

\section{Numerical Considerations and Edge Cases}

\begin{itemize}
    \item \textbf{Order of Reset Gate in Linear Layer:} Some implementations (e.g. cuDNN GRU) compute $U_h \mathbf{h}_{t-1}$ first and then multiply by $\mathbf{r}_t$: $\tilde{\mathbf{h}}_t = \tanh(W_h x_t + \mathbf{r}_t \odot (U_h \mathbf{h}_{t-1}) + \mathbf{b}_h)$, enabling fused GEMM across all three gates.
\end{itemize}

\section{References}
\begin{enumerate}
    \item Cho, K., Van Merri{\"e}nboer, B., Gulcehre, C., Bahdanau, D., Bougares, F., Schwenk, H., & Bengio, Y. (2014). Learning phrase representations using RNN encoder-decoder for statistical machine translation. \emph{EMNLP}, 1724--1734.
    \item Chung, J., Gulcehre, C., Cho, K., & Bengio, Y. (2014). Empirical evaluation of gated recurrent neural networks on sequence modeling. \emph{arXiv preprint arXiv:1412.3555}.
\end{enumerate}

\end{document}
